% Chapter 1 -- Electricity, circuits and the workbench
\chapter{Electricity, circuits and the workbench}
\label{ch:workbench}

This chapter builds the vocabulary and the bench habits that every later chapter
takes for granted. It answers four questions: what electricity \emph{is}
(as much as an engineer needs to know), what voltage, current and resistance
measure, how to read a schematic and turn it into a shopping list, and which of
the things on a workbench can hurt you. When you finish it you should be able
to: translate a simple schematic into parts and measurements, use a multimeter
for the three basic measurements without blowing its fuse, say precisely what
this book will and will not ask you to touch, and check whether your computer
already has the tools this book needs.

\section{What electricity is}
\label{sec:charge}

Strip everything away and electricity starts with one observation: objects can
push and pull on each other without touching. Rub a balloon on your hair and it
will pick up scraps of paper from a table. Something invisible reached out and
grabbed them. That something is \emph{charge}\index{charge|textbf}: a basic
property of matter, like mass, that some particles carry and others do not.
Charge comes in two flavours that attract each other and repel their own kind,
which we call positive and negative — the names are a bookkeeping device, not a
value judgment. The electron carries one fixed small negative charge; the proton
carries the same amount positive. That fixed amount, $e \approx
1.602\times10^{-19}$~coulombs\index{coulomb}, is the smallest packet of free
charge we ever deal with. Ordinary matter contains roughly equal numbers of
protons and electrons and so looks neutral; ``charging'' an object means
unbalancing that equality slightly.

A charge creates a disturbance in the space around it — an \emph{electric
field}\index{electric field|textbf} — and other charges feel that disturbance as
a force. We almost never draw the field itself in this book, but it is what
makes everything else work, so we define the quantity that packages it:
\emph{voltage}\index{voltage|textbf} (also called \emph{potential
difference}\index{potential difference}). Voltage is the field's push measured
as energy: the voltage between two points is the energy needed to move one unit
of charge from one to the other,
\[
  V = \frac{W}{q}, \qquad \text{1 volt} = \text{1 joule per coulomb}.
\]
A ``\qty{9}{\volt} battery'' is a device that gives each coulomb of charge
\qty{9}{\joule} of push between its terminals. Voltage is always measured
\emph{between two points}; the phrase ``the voltage at this wire'' is shorthand
for ``the voltage between this wire and whatever we agreed to call zero,'' and
that agreed-zero is the \emph{ground}\index{ground} node we draw at the bottom
of every schematic.

\begin{remark}[What we are skipping]
That paragraph is honest but thin. The full story is electromagnetic field
theory — fields propagate, wires are not the whole circuit, and at high
frequencies our tidy model breaks. For the circuits in this book (centimetres
across, signals slower than tens of megahertz) the \emph{lumped} model — every
property lives inside components connected by ideal wires — is accurate to far
better than the tolerance of any part you can buy. Where it stops being true
(radio, very fast digital edges, mains safety) we say so at the point it matters
rather than here. If you want the real derivations, a free calculus-based
physics text covers charge, field and potential properly \cite{openstaxv2}.
\end{remark}

\begin{figure}[htbp]
\centering
% Figure: charge -> moving charge -> current  (ch01)
\begin{tikzpicture}[scale=0.9,font=\small]
  % --- left: one atom ---
  \begin{scope}
    \fill[accentsoft] (0,0) circle (0.32);
    \node at (0,0) {\tiny +};
    \draw[slate!70] (0,0) circle (0.85);
    \draw[slate!70] (0,0) circle (1.25);
    \fill[ink] (0.85,0) circle (0.055) node[below right,xshift=1pt,yshift=-1pt]{\scriptsize $-$};
    \fill[ink] (-0.5,0.7) circle (0.055);
    \fill[ink] (-1.05,-0.66) circle (0.055);
    \node[align=center] at (0,-1.75) {\scriptsize matter: a few\\[-1pt]\scriptsize spare electrons};
  \end{scope}
  % --- arrow ---
  \draw[->,thick,accent] (1.7,0) -- (3.1,0);
  % --- middle: a wire full of drifting electrons ---
  \begin{scope}[xshift=3.4cm]
    \draw[rounded corners,thick] (0,-0.55) rectangle (4.6,0.55);
    \foreach \x/\y in {0.35/0.15, 0.8/-0.3, 1.25/0.35, 1.7/-0.1, 2.2/0.25, 2.7/-0.35, 3.15/0.1, 3.6/-0.2, 4.05/0.3, 4.35/-0.05}
      \fill[ink] (\x,\y) circle (0.05);
    \foreach \x in {0.5,1.5,2.5,3.5}
      \draw[->,slate] (\x,-0.75) -- (\x+0.5,-0.75);
    \node at (2.3,-1.35) {\scriptsize a wire: charge drifting along it};
    \node[slate] at (2.3,-1.05) {\scriptsize electron drift $\longrightarrow$};
  \end{scope}
  % --- right: the measure ---
  \begin{scope}[xshift=8.6cm]
    \draw[->,very thick,accent] (0,0) -- (1.6,0) node[right]{$I$};
    \node[align=center] at (0.8,-1.35) {\scriptsize current: how much\\[-1pt]\scriptsize charge per second};
  \end{scope}
\end{tikzpicture}

\caption{Read left to right: atoms are mostly neutral; a metal wire contains
electrons that are free to drift; a battery (the field source) makes them drift
in one direction. That drift \emph{is} the current. Notice what is \emph{not}
shown: the battery does not inject electrons into an empty pipe — the electrons
are already in the wire; the source only organises their motion.}
\label{fig:charge}
\end{figure}

\section{Current}
\label{sec:current}

A metal wire is full of electrons that belong to no particular atom and drift
randomly, as much left as right: no net motion, no current. Put a source across
the wire — a field along it — and the random drift gains a small organised
component in one direction. That net flow of charge is
\emph{current}\index{current|textbf}:
\[
  I = \frac{\Delta Q}{\Delta t}, \qquad \text{1 ampere} = \text{1 coulomb per second}.
  \index{ampere}
\]
One amp is roughly $6\times10^{18}$ electrons per second squeezing through a
cross-section of wire. You never count electrons; the meter does it for you by
watching the field they drag with them.

It is worth calibrating the unit with real numbers. The LED we light in this
chapter runs at about \qty{20}{\milli\ampere}. A phone charger supplies
\qtyrange{1}{2}{\ampere}. The LED strips in cheap room lighting draw a few
hundred milliamps. A bolt of lightning peaks in the tens of kiloamps — which is
why ``it is only a little circuit'' is a sentence that needs a voltage and a
current attached before it means anything.

Two conventions matter immediately. First, current is a \emph{signed}
quantity with an arrow: we pick a direction for the arrow, compute or measure
$I$, and a negative answer means the flow goes opposite to our arrow. Second —
and this trips up every beginner — the arrow points the way \emph{positive}
charge would move, which is opposite to the actual electron drift. That
historical quirk is Section~\ref{sec:conventions}'s whole topic; for now just
note that the arrow in Figure~\ref{fig:first} points the engineering way.

\section{Resistance and the resistor}
\label{sec:resistance}

Push charge through almost any material and it meets drag: the moving electrons
scatter off the lattice and hand their energy back as heat. The measure of that
drag is \emph{resistance}\index{resistance|textbf},
\[
  R = \frac{V}{I}, \qquad \text{1 ohm} = \text{1 volt per ampere}.
\]
A material with twice the resistance needs twice the voltage to push the same
current. Copper wire has nearly none (milliohms per metre); the carbon film
inside a quarter-watt \emph{resistor}\index{resistor|textbf} — the single most
common component in this book — is manufactured to a precise value from
fractions of an ohm to tens of megohms.

The relationship $V = I\,R$ is \emph{Ohm's law}\index{Ohm's law}: one equation
linking all three quantities, so that knowing any two fixes the third. We use
it lightly this chapter (to estimate a battery's life and to \emph{pose} the
problem in Example~\ref{ex:led}), and Chapter~2 turns it into a working tool.
\begin{remark}
``Ohm's law'' is really an empirical pattern, not a law of nature — it holds
beautifully for resistors and not at all for LEDs, transistors or arcs. The
components that break it are exactly the interesting ones; that is Chapters 5
and 6.
\end{remark}

Resistor values are printed as coloured bands\index{resistor!colour code} — a
holdover from when components were too small to label. You need only the
four-band reading: first two bands are digits, the third is the multiplier
(power of ten), the fourth is tolerance. Table~\ref{tab:bands} gives the digits;
Exercise~\ref{ex:easy-bands} drills it.

\begin{table}[htbp]
\centering\small
\begin{tabularx}{\linewidth}{lY}
\toprule
\thead{Colour} & \thead{Digit (multiplier $10^{n}$)}\\
\midrule
Black & 0\\
Brown & 1\\
Red & 2\\
Orange & 3\\
Yellow & 4\\
Green & 5\\
Blue & 6\\
Violet & 7\\
Grey & 8\\
White & 9\\
\midrule
Gold (4th band) & tolerance $\pm 5\%$\\
Silver (4th band) & tolerance $\pm 10\%$\\
\bottomrule
\end{tabularx}
\caption{The only colour-code facts this book needs: digit values and the two
common tolerance bands. Brown--black--red--gold $=1,0,\times10^{2},\pm5\% =
\qty{1}{\kilo\ohm}$. Notice you read \emph{toward} the gap between the digit
bands and the spaced-off tolerance band.}
\label{tab:bands}
\end{table}

Values span so many orders of magnitude that electronics runs on
\emph{engineering prefixes}\index{metric prefixes} in steps of $10^{3}$:
milli~m ($10^{-3}$), micro~µ (``u'' in plain ASCII, $10^{-6}$), nano~n
($10^{-9}$), pico~p ($10^{-12}$), kilo~k ($10^{3}$), mega~M ($10^{6}$). A
``4k7'' resistor is \qty{4.7}{\kilo\ohm} $= \qty{4700}{\ohm}$; a ``100 nF''
capacitor is \qty{100}{\nano\farad} $= \qty{0.1}{\micro\farad}$. Datasheets,
schematics and meters each pick different prefixes, so fluency is not optional.
The chapter ships a tiny converter (Listing~\ref{lst:prefix}) you can check your
mental arithmetic against.

\begin{exercise}[easy]
Convert: \qty{0.05}{\ampere} to milliamps; \qty{4.7}{\kilo\ohm} to ohms;
\qty{0.0000001}{\farad} to nanofarads. Then run \texttt{prefix.py} and check.
\end{exercise}

Before leaving the three quantities, one deserved warning about how you will
hear them explained elsewhere. The \emph{water analogy}\index{water analogy} —
voltage is pressure, current is flow, a resistor is a narrow pipe — is the
standard pictorial model, and it is not useless: it correctly predicts that
more pressure drives more flow and that pinching the pipe reduces it.
Figure~\ref{fig:analogy} puts it beside the honest version. It fails in
exactly the places where intuition matters most (a cut wire does not ``leak''
current onto the bench; a battery does not inject anything) — which is why
this book uses it as a first sketch and then abandons it.
Exercise~\ref{ex:water} asks you to catalogue the failures properly.

\begin{figure}[htbp]
\centering
% Figure: the water analogy, honestly framed (ch01)
\begin{tikzpicture}[scale=0.88,font=\small]
  % ---------- left: water ----------
  \begin{scope}
    % pump
    \draw[thick] (0,0) circle (0.5);
    \node at (0,0) {\scriptsize pump};
    % pipe loop
    \draw[very thick,accent] (0,0.5) -- (0,1.3) -- (4.2,1.3) -- (4.2,0.2);
    \draw[very thick,accent] (0,-0.5) -- (0,-1.3) -- (4.2,-1.3) -- (4.2,-0.2);
    % constriction
    \draw[very thick,accent] (4.2,0.2) -- (4.2,-0.2);
    \draw[thick,fill=graybg] (4.2,0.2) -- (4.55,0.08) -- (4.55,-0.08) -- (4.2,-0.2) -- cycle;
    % flow arrows
    \draw[->,thick] (1.4,1.05) -- (2.2,1.05);
    \draw[->,thick] (2.8,-1.05) -- (2.0,-1.05);
    \node[align=left,anchor=west] at (5.0,0.55) {\scriptsize pump $\approx$ source};
    \node[align=left,anchor=west] at (5.0,0.1) {\scriptsize pressure $\approx$ voltage};
    \node[align=left,anchor=west] at (5.0,-0.35) {\scriptsize constriction $\approx$ resistor};
    \node[align=left,anchor=west] at (5.0,-0.8) {\scriptsize flow $\approx$ current};
    \node at (2.3,-2.0) {\scriptsize\itshape the analogy};
  \end{scope}
  % ---------- right: where it breaks ----------
  \begin{scope}[yshift=-3.4cm]
    \node[align=left,anchor=north west,text width=13.2cm] at (-0.4,0)
      {\scriptsize\textbf{Where it breaks:} cut a pipe and the water pours out; cut a wire and
       charge simply stops. There is no ``leak'' of charge, nothing pushes sideways, and the
       pressure exists \emph{between} two points, not inside one point. Keep the picture for
       intuition in this chapter; drop it the moment it argues back.};
  \end{scope}
\end{tikzpicture}

\caption{The water analogy, kept honest. Left: the picture as usually drawn
(pump = battery, pipe = wire, constriction = resistor). Right: what each piece
actually is — pressure is energy per charge, flow is charge per second, the
``constriction'' is energy lost to lattice collisions. The analogy survives
Ohm's-law arithmetic and dies at the physics: a cut pipe leaks, a cut wire
just stops.}
\label{fig:analogy}
\end{figure}

\section{A complete circuit — and the schematic as the real object}
\label{sec:schematic}

Current flows only in loops. Three things must exist for it to do so: a
\emph{source} that provides the push (a battery, a USB port, a power supply), a
\emph{conducting path} that carries the flow (wire, plus the components in the
way), and a route \emph{back to the source}\index{circuit!closed loop}. Break
the loop anywhere and the current stops everywhere — instantly, throughout the
whole loop. That is the entire operating principle of a switch, and it is the
reason one broken wire can kill a whole circuit.

Figure~\ref{fig:first} shows the smallest circuit that does something useful:
a \qty{5}{\volt} source, a resistor, and a light-emitting diode (LED). The
schematic\index{schematic|textbf} — the drawn map of symbols and wires — is not
an illustration of the circuit; in every practical sense it \emph{is} the
circuit. The copper and plastic are just a heavy, solderable printout of the
schematic. This is why we insist all circuits in this book be drawn properly
(with \texttt{circuitikz}) rather than photographed.

\begin{figure}[htbp]
\centering
% Figure: first real schematic -- USB 5V, resistor, LED (ch01)
\begin{circuitikz}[scale=1.0,transform shape]
  \draw (0,0) node[ground]{};
  \draw (0,0) to[battery1,l_=\qty{5}{\volt},invert] (0,2.4);
  \draw (0,2.4) to[R,l=$R$] (3.2,2.4);
  \draw (3.2,2.4) to[leD,l=LED] (3.2,0) node[ground]{};
  % current arrow on top wire
  \draw[->,thick,accent] (0.6,2.85) -- (1.8,2.85) node[midway,above]{$I$};
  % node dots (the two wires between parts are "the same node")
  \node[circ] at (0,2.4) {};
  \node[circ] at (3.2,0) {};
\end{circuitikz}

\caption{Our first real schematic. Read it as a loop, not a picture: leave the
top of the source at $+\qty{5}{\volt}$, pass through $R$ (which will limit the
current), through the LED (which lights if the current is right), down to the
ground symbol and back into the bottom of the source. The arrow marks the
\emph{chosen} current direction (Section~\ref{sec:conventions}). Notice the LED
symbol: diode triangle plus two small arrows = ``the diode that emits.''}
\label{fig:first}
\end{figure}

Here is the subtle point beginners miss. The schematic records
\emph{connections}, not \emph{locations}. Two points joined by a drawn wire (or
marked with the same net name) are the same electrical object — the same
\emph{node}\index{node|textbf}. Points merely near each other on paper are
nothing. Figure~\ref{fig:topology} draws the same circuit twice: tidy, and
stretched into an ugly pose. Electrically they are identical, because every
node contains the same component leads. Reading schematics as
\emph{topology}\index{topology} — who is connected to whom, never where they
sit — is the skill this section exists to install. We will reuse it
mercilessly when schematics get big in Chapter~17.

\begin{figure}[htbp]
\centering
% Figure: same topology, different geography (ch01)
\begin{circuitikz}[scale=0.9,transform shape]
  % (a) canonical loop
  \begin{scope}
    \draw (0,0) node[ground]{};
    \draw (0,0) to[battery1,l_=\qty{5}{\volt},invert] (0,2.2);
    \draw (0,2.2) to[R,l=$R$] (2.8,2.2);
    \draw (2.8,2.2) to[leD,l=LED] (2.8,0) node[ground]{};
    \node at (1.4,-0.9) {\scriptsize (a) the tidy drawing};
  \end{scope}
  % (b) same netlist, bent around
  \begin{scope}[xshift=5.6cm]
    \draw (0,0) node[ground]{};
    \draw (0,0) to[battery1,l_=\qty{5}{\volt},invert] (0,2.2);
    \draw (0,2.2) -- (1.6,2.2) -- (1.6,0.9) to[R,l=$R$] (3.4,0.9);
    \draw (3.4,0.9) -- (3.4,1.8) to[leD,l=LED] (4.9,1.8) -- (4.9,0) node[ground]{};
    \node at (2.45,-0.9) {\scriptsize (b) bent, stretched --- still the same circuit};
  \end{scope}
\end{circuitikz}

\caption{Same circuit, two drawings. To prove (a) $\equiv$ (b), list what is
attached to each node: node A holds source-$+$ and resistor lead 1; node B holds
resistor lead 2 and LED anode; node C holds LED cathode and source$-$. The lists
match, so the circuits are the same. The drawing changed; the
\emph{connections} — the only thing the circuit cares about — did not.}
\label{fig:topology}
\end{figure}

\begin{example}[From schematic to shopping list]
\label{ex:partslist}
\emph{Problem.} Turn Figure~\ref{fig:first} into a parts order.\par
\emph{Strategy.} Walk the schematic component by component, not wire by wire.
For each symbol write the thing to buy and the value or rating written beside
it. Wires and ground symbols buy nothing — they come free with the breadboard
and cables.\par
\emph{Solution.}
\begin{itemize}
\item one \qty{5}{\volt} source — a USB cable or a \qty{5}{\volt} supply;
\item one resistor, value $R$ — value still undetermined
(Example~\ref{ex:led} is about exactly this gap);
\item one LED, any colour, standard \qty{5}{\milli\metre} or
\qty{3}{\milli\metre} package;
\item means of connecting: breadboard + jumper wires (Section~\ref{sec:kit}).
\end{itemize}
\emph{Notice:} the schematic already told us everything except one number —
which is the kind of gap a circuit diagram is supposed to leave for analysis.
\end{example}

\begin{example}[The LED question, posed on purpose]
\label{ex:led}
\emph{Problem.} The LED in Figure~\ref{fig:first} wants about \qty{2}{\volt}
across itself and shines nicely at \qty{10}{\milli\ampere}, burns at
\qty{30}{\milli\ampere}. The source is \qty{5}{\volt}. What should $R$
be?\par
\emph{Strategy.} Take stock of what each quantity does. The source fixes the
total push: \qty{5}{\volt}. The LED spends about \qty{2}{\volt} of it (a fact of
diode physics we take on trust until Chapter~5). The resistor must drop the
rest, and the same current must pass through everything in a single loop. So
the missing ingredient is a law that links the resistor's voltage share to the
current through it — which is $V = IR$. We \emph{have} that law already; what
we are missing is practice using it and one bookkeeping rule (Kirchhoff's
voltage law) that makes ``the rest'' precise.\par
\emph{Solution.} Held for Chapter~2, where it becomes a one-line computation.
The point of posing it now: you can already see that $R$ exists to dispose of a
voltage difference and to choose a current — the two jobs resistors do all book
long.
\end{example}

\section{Two conventions for direction}
\label{sec:conventions}

Benjamin Franklin guessed the sign of charge wrongly, and by the time electrons
were identified the convention was set: circuit arrows point the way
\emph{positive} charge would flow — \emph{conventional
current}\index{conventional current|textbf}. Electrons, being negative, drift
the other way — \emph{electron flow}\index{electron flow}.

Nothing about this is optional to get right, because everyone — every
datasheet, every schematic, every simulator — uses conventional current. When
someone (or a YouTube video) insists on electron flow, mentally flip every
arrow; the numbers and the physics are unaffected. The failure mode is not
choosing the ``wrong'' convention but switching conventions
\emph{mid-argument}: if you label an arrow one way and reason about electrons
the other way, signs invert silently and the sums stop making sense. Pick
conventional current and stay there; this book never leaves it.

\begin{principle}[Arrows are reference directions]
Draw the arrow whichever way is convenient. A positive computed value means
flow along your arrow; negative means against it. You do not have to guess the
``true'' direction first — the sign will tell you.
\end{principle}

\section{The multimeter}
\label{sec:dmm}

The \emph{multimeter}\index{multimeter|textbf} (or DMM, for digital
multimeter) is the one instrument this book treats as essential. It measures
the three quantities of Sections~1--3 and nothing else we need. Everything else
in electronics is learned through it. Figure~\ref{fig:dmm} shows a typical
unit.

\begin{figure}[htbp]
\centering
% Figure: a multimeter -- dial and probe jacks (ch01)
\begin{tikzpicture}[scale=0.95,font=\small]
  % body
  \draw[thick,rounded corners=8pt] (-2.3,-3.4) rectangle (2.3,3.4);
  % screen
  \draw[thick,fill=graybg] (-1.4,2.35) rectangle (1.4,3.05);
  \node at (0,2.7) {\scriptsize\texttt{1.990}};
  % dial
  \draw[thick] (0,0.35) circle (1.5);
  \fill[accentsoft] (0,0.35) circle (1.18);
  % needle
  \draw[thick,accent,->] (0,0.35) -- (210:1.05);
  \fill[accent] (0,0.35) circle (0.09);
  % dial markings around
  \foreach \a/\t in {150/{V$\sim$}, 120/{V\textbf{---}}, 90/{$\Omega$}, 60/{mA}, 30/{10\,A}, 0/{OFF}, -30/{cont.}, -60/{CAP}}
    \node[font=\scriptsize] at (\a:1.95) {\t};
  % jacks
  \draw[thick,fill=white] (-1.3,-2.3) circle (0.32) node[below,yshift=-7pt]{\scriptsize 10\,A};
  \draw[thick,fill=white] (0,-2.3) circle (0.32) node[below,yshift=-7pt]{\scriptsize COM};
  \draw[thick,fill=white] (1.3,-2.3) circle (0.32) node[below,yshift=-7pt]{\scriptsize V$\Omega$};
  % probes
  \draw[very thick] (-1.6,-2.9) -- (-2.5,-4.3) node[pos=1.05,below]{\scriptsize red};
  \draw[very thick] (0,-2.62) -- (-0.4,-4.4) node[pos=1.05,below]{\scriptsize black};
  \node[slate] at (0,-5.1) {\scriptsize red into V$\Omega$ for volts and ohms; the 10\,A jack is for current only};
\end{tikzpicture}

\caption{A typical multimeter. Three things to notice. (1)~The dial divides the
world into ranges: $\mathrm{V}\!\sim$ (AC volts — never used in this book),
$\mathrm{V}\!\text{-}\!\!\cdot$ (DC volts), $\Omega$ (resistance), mA (current),
10\,A (the big current range), plus extras (continuity beeper, capacitance).
(2)~Three jacks: COM is where the black probe always lives; V$\Omega$ is red for
volts, ohms and small currents; 10\,A is a separate, typically \emph{unfused}
path for large currents. (3)~The screen reads the number \emph{in the selected range's units} —
``1.990'' on the 2\,V range means \qty{1.990}{\volt}, on the mA range
\qty{1.990}{\milli\ampere}.}
\label{fig:dmm}
\end{figure}

The three measurements differ in a way that matters physically, and the
difference explains every DMM accident:

\begin{itemize}
\item \textbf{Voltage} is a \emph{difference between two points}. To measure
it you touch the two probes to the two points — the meter is in
\emph{parallel}\index{multimeter!measuring voltage} with whatever is there.
The circuit does not change; you added an observer, not a wire. Inside the
meter is a very large resistance ($\sim\qty{10}{\mega\ohm}$), so almost no
current diverts into it. Voltage measurement is the safe, always-allowed move.
\item \textbf{Current} is a \emph{flow through a point}. To measure it you must
put the meter \emph{into} the path — break the circuit\index{multimeter!measuring current}
and bridge the gap with the probes, so all the circuit's current passes through
the meter. Inside it is a very \emph{small} resistance (a shunt). This is the
measurement that changes the circuit and the one beginners fry meters with.
\item \textbf{Resistance} is measured with the circuit \emph{unpowered}: the
meter sends out its own small test current and reads the resulting voltage.
Measuring resistance on a live circuit gives wrong numbers at best.
\end{itemize}

\begin{tripwire}[The fuse-blower]
Red probe in the current jack + probes across a voltage source = a short
circuit through the meter's (very small) input resistance. On most meters the
mA jack is the fused one — a sacrificial fuse dies for your mistake; the
10\,A jack is typically \emph{unfused}, so the same mistake there can kill
the meter itself, not to mention the circuit you shorted. Make a habit: \emph{after} a
current measurement, the red probe goes back to V$\Omega$ before anything else
is touched. We will never again let it live in the current jack when idle.
\end{tripwire}

\begin{tripwire}[``OL'' is an answer]
On a resistance or voltage range, an ``OL''/``1\ \ '' display means
\emph{over-range}, not zero and not broken: the value exceeds what that range
can show (or the path is genuinely open). On the mA current jack by
contrast, ``OL'' with a live circuit almost always means the range's fuse
blew — that fuse exists precisely for this accident. Same screen, opposite
meanings — know which.
\end{tripwire}

\begin{checkpoint}[Predict before you measure]
A fresh 9\,V battery reads $\sim$9.4--9.6\,V on a meter, not 9.000. An
``empty'' one reads $\sim$7\,V. \emph{Before} every measurement this book asks
you to make, write down your expected number and its range. If the meter
disagrees wildly, the mistake is in the setup, not the meter — check the dial
and the jacks first.
\end{checkpoint}

\begin{exercise}[medium]
You want to measure the current through the LED of
Figure~\ref{fig:first}. State exactly what you physically change in the
circuit, which range and jacks you use, and where each probe tip lands.
\end{exercise}

\section{The workbench kit and the budget}
\label{sec:kit}

This book is built to need almost nothing bought. Every analogue circuit is
simulated in ngspice before (or instead of) being built; the entire digital
half runs in simulation on free tools; the one hard purchase — a Tang Nano 9K
FPGA board — is needed only from Chapter~13 and costs about US\$20. But a
minimum physical kit makes several chapters concrete, so Table~\ref{tab:kit}
lists it with honest prices.

\begin{table}[htbp]
\centering\small
\begin{tabularx}{\linewidth}{lYl}
\toprule
\thead{Item} & \thead{What it is for} & \thead{Budget (USD)}\\
\midrule
Multimeter & every chapter's measurements; buy auto-ranging, fused & \$15--30\\
Breadboard + jumpers & solderless builds, ch17--19 & \$5--8\\
Resistor assortment (E24) & every analogue build & \$4--6\\
LEDs, a few colours & ch2, ch5, board bring-up & \$2--4\\
USB cable / 5\,V supply & power for bench circuits & \$3--5\\
9\,V battery + clip & the battery-life experiment below & \$2--3\\
\midrule
\textbf{Physical kit subtotal} & & \textbf{\$31--56}\\
Tang Nano 9K (Sipeed, GW1NR-LV9) & FPGA chapters 13--16, required then & \${18--23}~\cite{tangnano9k}\\
\bottomrule
\end{tabularx}
\caption{The shopping list as a figure: what a minimal bench costs in 2026.
Notice the split — the top block is a \emph{recommended} bench; the FPGA board
is the only \emph{required} purchase, and only from Chapter~13. Indonesian
availability/price for the board is marked unverified in the sources log; the
global street price is US\$18--23 \cite{tangnano9k}.}
\label{tab:kit}
\end{table}

Two things about this list deserve saying. First, a simulator never burns
fingers and never needs resoldering — every analogue circuit in this book comes
with an ngspice run you can do for free, and the experiments are structured so
hardware is always a checkpoint, not a gate. Second, the board matters later:
the custom CPU you design in Chapters 14--16 runs on real silicon, and the
difference between ``simulated'' and ``it blinks on a board'' is the whole
reason Part~V exists.

The tools, unlike the kit, must be checked on your computer now.
Listing~\ref{lst:checktools} probes each command the book needs and reports
what is installed; Listing~\ref{lst:checktools-out} is its real output on a
fresh machine — read the MISSING lines as a to-install list, not an error.

\begin{lstlisting}[language=bash,caption={\texttt{code/check\_tools.sh} --
reports which toolchain commands exist on your PATH.},label={lst:checktools}]
#!/bin/sh
# check_tools.sh -- verify the toolchain for this book is installed.
# Run: sh code/check_tools.sh
# Goal: every line says "found". Anything MISSING tells you what to install.

check() {
  # $1 = command to try, $2 = version flag, $3 = what it is for
  if command -v "$1" >/dev/null 2>&1; then
    ver=$("$1" $2 2>/dev/null | head -n 1)
    printf 'found    %-22s %s\n' "$1" "$ver"
  else
    printf 'MISSING  %-22s %s\n' "$1" "$3"
  fi
}

check iverilog    -V        "Verilog simulation (ch11+)"
check verilator   --version "lint / fast simulation (ch12+)"
check yosys       -V        "synthesis (ch13+)"
check nextpnr-himbaechel --version "place-and-route, Gowin (ch13+)"
check gowin_pack  --help    "bitstream packer, Gowin (ch13+)"
check openFPGALoader --Version  "programming the board (ch13+)"
check gtkwave     --version "viewing waveforms (ch11+)"
check ngspice     --version "analog simulation (ch2+)"
check kicad-cli   version   "PCB work (ch18)"
check python3     --version "scripts and vector generators"

echo "----"
echo "Missing tools come from the OSS CAD Suite (FPGA tools) and your"
echo "package manager (ngspice, KiCad). See code/README.md for versions."
\end{lstlisting}

\begin{lstlisting}[language={},numbers=none,caption={Real output of
\texttt{check\_tools.sh} on a machine before installing the toolchain. Every
MISSING line names both the command and where it will first be needed.},
label={lst:checktools-out}]
MISSING  iverilog               Verilog simulation (ch11+)
MISSING  verilator              lint / fast simulation (ch12+)
MISSING  yosys                  synthesis (ch13+)
MISSING  nextpnr-himbaechel     place-and-route, Gowin (ch13+)
MISSING  gowin_pack             bitstream packer, Gowin (ch13+)
MISSING  openFPGALoader         programming the board (ch13+)
MISSING  gtkwave                viewing waveforms (ch11+)
MISSING  ngspice                analog simulation (ch2+)
MISSING  kicad-cli              PCB work (ch18)
found    python3                Python 3.9.6
----
Missing tools come from the OSS CAD Suite (FPGA tools) and your
package manager (ngspice, KiCad). See code/README.md for versions.
\end{lstlisting}

The second small program converts between the engineering prefixes that meters
and datasheets disagree about (Listing~\ref{lst:prefix}); its output is
Listing~\ref{lst:prefix-out}. Neither program is mandatory — both exist so the
chapter has real, runnable code like every other one.

\begin{lstlisting}[language=Python,caption={\texttt{code/prefix.py} --
metric-prefix helper used in the exercises.},label={lst:prefix}]
#!/usr/bin/env python3
"""prefix.py -- metric-prefix converter used in Chapter 1 exercises.

Given a value like 0.0047 it prints engineering notation (4.7 m).
Why: datasheets and meters never agree on the prefix; being able to
shift between them fluently is the first practical skill.
"""

PREFIXES = [
    (1e9, "G"), (1e6, "M"), (1e3, "k"),
    (1.0, ""), (1e-3, "m"), (1e-6, "u"), (1e-9, "n"), (1e-12, "p"),
]

def to_eng(value):
    """Format a plain number in engineering (prefix-aligned) notation."""
    if value == 0:
        return "0"
    for factor, sym in PREFIXES:
        if abs(value) >= factor:
            return f"{value / factor:g} {sym}".rstrip()
    return f"{value:g}"

if __name__ == "__main__":
    for v in [4700, 0.0047, 0.0000001, 20e-3, 3.3]:
        print(f"{v:g}  ->  {to_eng(v)}")
\end{lstlisting}

\begin{lstlisting}[language={},numbers=none,caption={Real output of
\texttt{python3 code/prefix.py}. Compare the third line: $10^{-7}$ F is
the ceramic-disc capacitor marked ``104'' you will meet in Chapter~4.},
label={lst:prefix-out}]
4700  ->  4.7 k
0.0047  ->  4.7 m
1e-07  ->  100 n
0.02  ->  20 m
3.3  ->  3.3
\end{lstlisting}

\begin{exercise}[medium]
Prove Figure~\ref{fig:topology}(a) $\equiv$ (b) formally: write the contents
of each node in each drawing and show the two tables are identical. (This is
Exercise~\ref{ex:easy-bands}'s harder sibling: the colour-code drill is
mechanical; this one tests whether the \emph{idea} of a node landed.)
\end{exercise}

\section{Worked example: how long does the battery last?}
\label{sec:battery}

We finish the technical content with a worked calculation of the sort this
book does constantly — a back-of-envelope estimate that turns component numbers
into an engineering decision.

\begin{example}[Battery life for a 20 mA LED]
\label{ex:battery}
\emph{Problem.} You power the LED of Figure~\ref{fig:first} from a 9\,V
battery instead of USB. How long will it run?\par
\emph{Strategy.} Batteries store charge, rated in
milliamp-hours\index{milliamp-hour}: the product ``mA $\times$ hours'' the
battery can deliver before it is flat. The naive estimate is capacity
$\div$ current. Then derate: real batteries deliver less than nominal,
especially at higher drain, so multiply by a real-world factor.\par
\emph{Solution.} A small 9\,V alkaline is rated about 550~mAh.
\[
  t \approx \frac{550\ \text{mAh}}{20\ \text{mA}} \times 0.7
    \approx 19\ \text{hours}.
\]
The $0.7$ is honesty, not precision — datasheet curves for this drain sit near
it. So: a day of continuous light, a week of occasional breadboard use.
\emph{Notice:} the answer came from one multiplication; the skill was knowing
that mAh is a charge reservoir and that the LED current, not the voltage, sets
the drain rate. The voltage instead decides \emph{how many components} the
battery can push through (Chapter~2's job).
\end{example}

\section{Safety: what this book will and will not have you touch}
\label{sec:safety}

Safety in a book like this is not a disclaimer page; it is a set of numbered
rules you can recite. Here they are, all of them, for the whole book.

\begin{rulebox}[The low-voltage rule]
Everything this book asks you to build runs at or below
\qty{9}{\volt} from batteries or a USB port. At these voltages, dry skin stops
essentially all current — you cannot get a dangerous shock from the circuits
in this book. What remains dangerous is \emph{heat}: a shorted battery, a
resistor asked to drop too much power, a reversed electrolytic capacitor. The
rule: \textbf{low voltage only, and if anything gets hot, disconnect power
first, then look.}
\end{rulebox}

What about the numbers that \emph{are} dangerous? Mains in Indonesia is
\qty{230}{\volt} at \qty{50}{\hertz}, fused at 10--16~A in the wall. Currents
through the body stop being safe somewhere around 10--20~mA sustained (the
exact threshold depends on path and person); mains can drive far more through
muscle and heart. You do not need drama, just arithmetic: $230$~V across the
$\sim$1--10~k$\Omega$ of skin gives tens to hundreds of milliamps — over the
threshold by an order of magnitude. The book's treatment of mains is therefore
complete in one rule.

\begin{rulebox}[The wall rule]
If it plugs into the wall, we do not open it. Not the charger brick, not the
supply inside it, not anything. This book never requires you to work on
mains voltage; do not work on it without proper supervision and equipment.
\end{rulebox}

Three more standing hazards, all quiet ones:

\begin{itemize}
\item \textbf{ESD} (electrostatic discharge\index{ESD}): the spark you feel
touching a doorknob is thousands of volts at trivial energy — harmless to you,
fatal to the MOSFETs and the FPGA inside Chapters 6 onward. Handle boards by
their edges; touch a grounded metal object first; keep ICs in their foam or
bags until use.
\item \textbf{Battery shorts}\index{battery!short}: an AA or 9\,V shorted
across its own terminals heats fast enough to burn. A lithium (LiPo) cell
shorted or punctured can enter thermal runaway — fire, not warmth. Never let
loose batteries share a pocket with keys or wire ends; never leave a LiPo
charging unattended.
\item \textbf{Hot solder and hot parts}: from Chapter 17 we solder. The iron
lives in its stand; parts get 30~seconds before fingers. A dropped hot iron is
caught by reflex — break the reflex now: let it fall.
\end{itemize}

\begin{checkpoint}[Rules audit]
Say these back without looking: the voltage ceiling for this book's builds;
the one place the red probe may not rest; what to do when a component is hot;
the four words of the wall rule. If any answer took more than a second, reread
this section — it is the only one where the stakes are not academic.
\end{checkpoint}

\section{Common mistakes}
\label{sec:mistakes01}

Each chapter ends the main text by naming the errors that actually hurt people.
Four belong to this chapter.

\begin{description}
\item[Reading the schematic as geography.] Positions on paper mean nothing;
connections mean everything. Two drawn wires that cross without a junction dot
do not connect; two points at opposite corners of the page joined by a net name
are the same node. Re-doing Figure~\ref{fig:topology}'s node-listing exercise
until it is automatic is the fix. \index{schematic!as topology}
\item[Mixing current conventions mid-argument.] Choosing an arrow, then
reasoning in electron flow, silently flips signs and produces $I = -20$~mA
``errors'' that are pure bookkeeping noise. Conventional current, always.
\index{conventional current}
\item[Thinking the meter ``uses up'' what it measures.] A voltmeter adds a
10~M$\Omega$ observer in parallel — it diverts essentially nothing. An
\emph{ammeter}, inserted in series, is the opposite: a near-short in the path.
The confusion between the two is exactly the fuse-blower tripwire of
Section~\ref{sec:dmm}. \index{multimeter}
\item[Forgetting that a wire is a component too.] Ideal schematics hide
resistance in wire and contacts; breadboards add ohms where fingers expect
zero. When a measured number disagrees with a simulation in later chapters,
the first suspect is the \emph{connection}, not the math.
\end{description}

\section{Chapter summary}
\label{sec:ch1summary}

\begin{table}[htbp]
\centering\small
\begin{tabularx}{\linewidth}{llY}
\toprule
\thead{Quantity / idea} & \thead{Symbol \& unit} & \thead{The one-line version}\\
\midrule
Charge & $q$,\,C & property of matter; electrons carry $-e$\\
Voltage & $V$,\,V & energy per charge, between two points ($V=W/q$)\\
Current & $I$,\,A & charge per second through a point ($I=\Delta Q/\Delta t$)\\
Resistance & $R$,\,$\Omega$ & volts needed per amp pushed ($R=V/I$)\\
Node & — & all points joined by wire/name = one object\\
Conventional current & — & arrows point positive flow; electrons go the other way\\
Multimeter & V, A, $\Omega$ & parallel for V; series (break it) for A; power off for $\Omega$\\
Metric prefixes & m, µ, n, p, k, M & $10^{-3}$ to $10^{6}$ in steps of $10^{3}$\\
Battery capacity & mAh & charge reservoir; life $\approx$ capacity/current $\times$ 0.7\\
\midrule
Safety & — & $\le 9$\,V builds only; never open wall-powered gear; ESD kills chips\\
\bottomrule
\end{tabularx}
\caption{Chapter 1 on a page. If a term in the left column does not produce the
right column instantly, the section named in its row is the one to reread.}
\label{tab:ch1summary}
\end{table}

\section*{Exercises}
\addcontentsline{toc}{section}{Exercises}
\label{sec:ch1exercises}

The exercises marked in the text were checkpoints; these stand alone.
[easy] is one step, [medium] a small argument, [hard] an open problem.

\begin{exercise}[easy]
\label{ex:units2}
Convert each to engineering notation, then verify with \texttt{prefix.py}:
\qty{0.0022}{\ampere}; \qty{10000}{\ohm}; \qty{0.000000047}{\farad}.
\end{exercise}

\begin{exercise}[easy]
\label{ex:easy-bands}
Using Table~\ref{tab:bands}, decode: (a) red--red--brown--gold;
(b) yellow--violet--orange--gold; (c) brown--black--green--gold.
\end{exercise}

\begin{exercise}[easy]
\label{ex:dmm-select}
For each measurement, name the dial position \emph{and} the two jacks:
(a) voltage of a 9\,V battery; (b) current through the LED in
Figure~\ref{fig:first}; (c) resistance of an unmounted \qty{1}{\kilo\ohm}
resistor.
\end{exercise}

\begin{exercise}[easy]
\label{ex:nodes}
How many distinct nodes does Figure~\ref{fig:first} contain? Name them.
(Count the ground node; it is one node.)
\end{exercise}

\begin{exercise}[medium]
\label{ex:samecircuit}
Two schematics are described in words. A: ``source, then LED, then resistor,
then ground.'' B: ``source, then resistor, then LED, then ground.'' Same
circuit or different? Justify with node lists — this has a correct
electrical answer, not a taste answer.
\end{exercise}

\begin{exercise}[medium]
\label{ex:aa-life}
Two AA alkaline cells in series make a nominal 3\,V, $\sim$2400~mAh pack.
Recompute Example~\ref{ex:battery} for a 15~mA draw, including the 0.7
derating. Then explain \emph{why} the derating exists — what is the battery
losing to? (Hint: a battery is not an ideal voltage source; it has
resistance inside it, called internal resistance, and real capacity falls as
drain rate rises.)
\end{exercise}

\begin{exercise}[medium — derive]
\label{ex:ledmissing}
An LED datasheet says: ``$V_F = 2.1$ V (typical), $I_F$ max = 30 mA.''
To pick $R$ for the circuit of Figure~\ref{fig:first} with a 5\,V source and a
target of 15~mA, list each quantity Ohm's law needs, state which of them the
datasheet gives, and write the one equation you are still a step short of
solving. (You are deriving the shopping list of Chapter~2, not solving it.)
\end{exercise}

\begin{exercise}[medium]
\label{ex:ol}
Your meter on the 10\,A current setting shows ``OL'' for every circuit you
try. Give two explanations consistent with the display — one in the circuit,
one in the meter — and a test that separates them.
\end{exercise}

\begin{exercise}[hard — essay]
\label{ex:water}
The water analogy (voltage = pressure, current = flow, resistor = narrow
pipe) is everywhere, and it is wrong in instructive ways. Write one page:
identify \emph{two} concrete places it misleads (suggestions: what happens
when a pipe ``ends'' vs.\ a wire ends; what a battery actually does to
electrons already in the wire; what ``pressure between two points'' means for
a hanging wire with no return path), and for each give the correct lumped-model
statement instead. Marking rubric in the hints.
\end{exercise}

\begin{exercise}[hard — design]
\label{ex:kitdesign}
You have \$40 for the kit of Table~\ref{tab:kit}, and the breadboard and
multimeter are non-negotiable. Choose the rest and defend the two items you
left out: for each, name a chapter capability you sacrificed and the
simulation that partially covers it. (There are at least three defensible
answers; grade is on the justification matching the book's structure.)
\end{exercise}

\subsection*{Hints and selected answers}
\addcontentsline{toc}{subsection}{Hints and selected answers}

\begin{description}
\item[\ref{ex:units2}] 2.2~mA; 10~k$\Omega$; 47~nF.
\item[\ref{ex:easy-bands}] (a) 220~$\Omega$; (b) 47~k$\Omega$;
(c) 1~M$\Omega$; all $\pm$5\%.
\item[\ref{ex:dmm-select}] (a) DC~V range, COM + V$\Omega$, probes across the
battery (parallel). (b) mA or 10\,A range, COM + current jack, \emph{break}
the loop and bridge with probes (series). (c) $\Omega$ range, COM +
V$\Omega$, resistor unpowered.
\item[\ref{ex:nodes}] Three: source$+$/resistor; resistor/LED anode; LED
cathode/ground/source$-$.
\item[\ref{ex:samecircuit}] Same circuit. In a single loop, element order does
not enter the node table — every element sits between the same two neighbours'
nodes. Node lists identical $\Rightarrow$ identical.
\item[\ref{ex:aa-life}] $2400/15 \times 0.7 = 112$~h. The battery's internal
resistance converts part of its energy to heat inside itself, and effective
capacity shrinks as current rises — the mAh rating assumes a gentle drain.
\item[\ref{ex:ledmissing}] Needed: resistor voltage $V_R$ and target current
$I$. Given: $V_F \approx 2.1$ V, source 5\,V $\Rightarrow V_R = 5-2.1=2.9$ V;
$I = 15$~mA. Missing equation solved next chapter: $R = V_R/I$.
\item[\ref{ex:ol}] Either the circuit's current path is open (a break
somewhere) or the meter's mA-range fuse is blown — the fuse lives on the
current-measurement path, not the (unfused) 10\,A jack. Test: measure the
meter's own resistance between COM and the mA jack — near 0 $\Rightarrow$
path good, OL $\Rightarrow$ blown fuse.
\item[\ref{ex:water}] Rubric: each failure must (1) state the wrong
prediction the analogy makes, (2) name the lumped fact that replaces it,
(3) fit in two sentences. Examples that pass: ``pipe ends $\Rightarrow$ leak''
vs.\ open circuit $\Rightarrow$ current simply stops; ``pump adds water'' vs.\
source reorganises drift of electrons already present.
\item[\ref{ex:kitdesign}] Strong answers drop, e.g., the 9\,V battery (the
USB supply covers every experiment) and spare LEDs (one suffices for ch2),
and note ngspice covers the battery-life check in simulation.
\end{description}

\chapternotes

\begin{enumerate}
\item T.~R. Kuphaldt, \emph{Lessons in Electric Circuits}, Vol.~I ``DC'',
chapters 1--3 (static electricity, conductors/insulators/electron flow,
electric safety). The sequencing theory-then-safety is inherited from here;
the safety facts are restated in our terms. Free online
\cite{aacvol1}.
\item OpenStax, \emph{University Physics} Vol.~2 (Ling, Moebs, Sanny et
al.): the charge, field and potential definitions that
Section~\ref{sec:charge} compresses; the reference for readers who want the
skipped derivations \cite{openstaxv2}.
\item SparkFun Electronics, ``How to Use a Multimeter'': the dial/jack
anatomy Figure~\ref{fig:dmm} redraws, and the fuse/jack warning
\cite{sparkfunmultimeter}.
\item SparkFun Electronics, ``How to Use a Breadboard'': internal connection
topology used later in ch17; cited here because the ``rails are not joined''
trap is first mentioned in Section~\ref{sec:kit}
\cite{sparkfunbreadboard}.
\item The water analogy's honest limits are standard physics-teaching fare;
OpenStax vol.~2's treatment of potential is the correct statement
\cite{openstaxv2}. (Exercise~\ref{ex:water} asks you to produce it, not
quote it.)
\item Tang Nano 9K specifications and the open-source toolchain for it
(yosys, nextpnr-himbaechel, \texttt{gowin\_pack}, openFPGALoader) are
documented by Project Apicula; the street price is marked [unverified] per
the decisions log \cite{tangnano9k}.
\end{enumerate}
